% =====================================================================
%  CHAPTER 13: 生长发育调查评价与健康危险行为监测（对应大纲第十三章）
%  参考：大纲第13章
% =====================================================================
\chapter{生长发育调查评价与健康危险行为监测}
\setcounter{chapter}{13}
\chapterintro{
\textbf{掌握：}生长发育调查设计（横断面、监测、纵向、序列、干预性设计）及质量控制；生长发育水平评价方法（离差法、百分位数法、Z分法）及等级划分；发育年龄评价（形态年龄、齿龄、骨龄*）；健康危险行为分类。\newline
\textbf{熟悉：}生长速度评价、匀称度评价；健康危险行为监测内容与应用。
}

% =====================================================================
% 第一部分：生长发育调查
% =====================================================================
\section{生长发育调查}

\subsection{调查设计}

\begin{keybox}
\textbf{（一）横断面设计（Cross-sectional Design）}

在某一时间段内，选择特定地区、有代表性的不同年龄阶段儿童少年，针对特定指标进行一次性的群体大规模调查。是生长发育调查的主要方法。

\begin{itemize}[leftmargin=*]
    \item \imp{优点}：可在短期内获得大量数据，费用相对较低，适用范围广。
    \item \imp{要求}：应有周详的调查方案，选取指标少而精；采取适宜的抽样方案（如多阶段分层随机整群抽样），使对象有较强代表性；检测方法和技术规范应统一。
    \item \imp{局限性}：仅反映某一时间点的状况，无法揭示生长发育的动态变化和因果关系。
\end{itemize}

\textbf{（二）监测设计（Surveillance Design）}

对某地区、某群体的生长发育指标进行连续收集、整理、分析的过程，是了解儿童少年生长发育动态变化的重要手段。

\begin{enumerate}[leftmargin=*]
    \item \imp{明确监测目的}：界定监测需解决的问题。
    \item \imp{确定监测人群}：可在学校、社区、医院等不同基础上对学龄儿童少年进行人群监测（population-based surveillance）。
    \item \imp{稳定监测方法}：监测体系固定，指标简便易行。
    \item \imp{保证质量控制}：保障结果的可比性。
    \item \imp{结果反馈与应用}：及时制定和调整干预措施。
\end{enumerate}

\textbf{（三）纵向设计（Longitudinal Design）/ 随访设计（Follow-up Design）}

在一段时间如几年乃至终生对同一批儿童少年个体进行反复观察的研究设计。

\begin{itemize}[leftmargin=*]
    \item \imp{优点}：研究对象是同一组个体，能揭示早期经历（不良环境、营养不良等）与后来生长发育结果间的因果关系。
    \item \imp{缺点}：耗时长，花费大，样本易流失；可能存在"学习效应"。
\end{itemize}

\textbf{（四）序列设计（Sequence Design）}

对不同年龄组的儿童少年进行横向调查，间隔一定时间后对同一批儿童少年进行一次或多次重复调查，构成追踪性研究，将横断面调查和追踪调查相结合。

\begin{itemize}[leftmargin=*]
    \item \imp{优点}：克服追踪调查所需年限太长、样本易流失等缺点；节约时间和工作量。
    \item \imp{缺点}：仅具部分追踪性质，所获生长速度数据是近似的。
\end{itemize}

\textbf{（五）干预性设计（Interventional Design）}

采用随机对照试验（RCT）设计，将对象随机分为干预组和对照组，对干预组施加某种干预措施（如营养补充、健康教育、体育锻炼等），通过比较两组生长发育指标的变化来评价干预效果。

\begin{itemize}[leftmargin=*]
    \item \imp{特点}：能最有力地证明因果关系，是评价干预措施效果的"金标准"。
    \item \imp{要求}：严格随机分组、设立对照、盲法实施、伦理审批。
\end{itemize}
\end{keybox}

\subsection{质量控制}

\begin{keybox}
\textbf{（一）设计阶段质量控制}

\begin{enumerate}[leftmargin=*]
    \item \imp{制定详细调查方案}：明确调查目的、对象、方法、指标、时间、人员、经费等。
    \item \imp{伦理学审查}：向伦理委员会提交申请报告，获得批准后方可实施。
    \item \imp{知情同意}：取得研究对象（14岁以下需家长）的知情同意。
    \item \imp{预调查}：正式调查前进行小规模预调查，检验方案的可行性和仪器设备。
\end{enumerate}

\textbf{（二）实施阶段质量控制}

\begin{enumerate}[leftmargin=*]
    \item \imp{检测仪器和方法}：仪器应精确，方法应统一；追踪调查应始终使用同一方法和同种仪器。
    \item \imp{检测时间}：追踪调查时个体测量时间应相对固定；将同龄样本均匀分配在上午和下午；考虑季节对生长发育的影响，一般选择5--6月或9--10月。
    \item \imp{年龄计算}：我国按实足年龄（测试年月日 $-$ 出生年月日）计算，满7岁到差1天满8岁者均为7岁。
    \item \imp{技术培训}：严格培训测试者，统一操作规范，考核合格方可上岗。
    \item \imp{检测程序}：各项目按规定顺序流水作业，避免漏测；合理安排检测流程。
\end{enumerate}

\textbf{（三）质量控制检查}

\begin{enumerate}[leftmargin=*]
    \item \imp{人员检验}：定期对检测人员进行操作一致性检验。
    \item \imp{现场检验}：
    \begin{itemize}
        \item 逐一核对调查表，检查缺项、漏项。
        \item 记录准确，字迹清晰。
        \item 计算误差率（$P$）：若 $P > 5\%$，应及时分析原因并采取措施改进；若 $P > 20\%$，当日全部检测数据无效，重新测试。
    \end{itemize}
    \item \imp{逻辑检验}：检查数据间的逻辑关系是否合理（如身高与坐高的关系等）。
\end{enumerate}
\end{keybox}

\subsection{调查应用}

\begin{defbox}
生长发育调查的主要应用领域：

\begin{enumerate}[leftmargin=*]
    \item \imp{描述生长发育水平}：了解本地区儿童少年生长发育现状，为制定卫生标准提供依据。
    \item \imp{反映生长发育动态}：对同地区、同人群进行连续多次调查，追踪儿童少年生长发育动态变化，分析生长长期趋势（secular trend）。
    \item \imp{评价队列效应（Cohort Effect）}：不同出生队列因经历不同时期，其生理及心理健康发展状况可反映相应阶段的社会变化与时代特征。
    \item \imp{评价干预措施效果}：通过比较干预前后的生长发育指标变化，客观评价学校卫生干预措施的有效性。
    \item \imp{制定生长发育参考标准}：基于大样本调查数据，建立本地区、本民族的生长发育正常值或评价标准。
    \item \imp{为卫生政策制定提供依据}：调查结果为政府制定儿童少年健康相关政策和规划提供科学依据。
\end{enumerate}
\end{defbox}

% =====================================================================
% 第二部分：生长发育评价
% =====================================================================
\section{生长发育评价}

\subsection{评价目标与标准}

\subsection{发育水平评价}

\begin{keybox}
\textbf{（一）离差法（Deviation Method）}

\textbf{1. 评价原理：}用均值（$\bar{x}$）和标准差（$s$）描述样本调查值的分布，适用于正态分布资料。一般以 $\bar{x} \pm 2s$ 为正常范围。$\bar{x} \pm s$ 包括样本的68.3\%，$\bar{x} \pm 2s$ 包括95.4\%，$\bar{x} \pm 3s$ 包括99.7\%。

\textbf{2. 评价方法——五等级评价标准：}

\begin{table}[H]
\centering
\caption{离差法五等级划分}
\begin{tabularx}{\textwidth}{ll}
\hline
\textbf{等级} & \textbf{标准} \\
\hline
上等   & $> \bar{x} + 2s$ \\
中上等 & $\bar{x} + s \sim \bar{x} + 2s$ \\
中等   & $\bar{x} \pm s$ \\
中下等 & $\bar{x} - 2s \sim \bar{x} - s$ \\
下等   & $< \bar{x} - 2s$ \\
\hline
\end{tabularx}
\end{table}

\textbf{3. 优缺点：}

\begin{itemize}[leftmargin=*]
    \item \imp{优点}：方便快捷，适用于基层；结果含义明确。
    \item \imp{缺点}：①只适用于正态分布指标，对偏态分布指标（如体重、BMI）评价不够精确；②结果只能针对单项指标，无法反映个体发育匀称程度；③不能直观反映动态变化。
\end{itemize}

\textbf{4. 注意事项：}个体实测值在 $\bar{x} \pm 2s$ 内视为正常；在 $\bar{x} \pm 2s$ 外不能轻易判定为异常，需连续观察，结合临床检查，慎重下结论。

\textbf{（2）曲线图法}

将同性别、各年龄组的均值、均值±1s、均值±2s分别连成曲线，制成生长发育标准曲线图。评价时将个体实测值直接点在曲线图上，根据点在图中的位置确定发育等级。

\begin{itemize}[leftmargin=*]
    \item \imp{优点}：使用简便、结果直观，可同时针对个体和群体进行评价；通过对同一年龄组不同年代的发育均值比较，可分析生长发育的长期变化趋势。
    \item \imp{缺点}：需事先分男女、针对不同指标制作相应曲线图；对非正态分布指标不够精确。
\end{itemize}

\textbf{（二）百分位数法（Percentile Method）}

\textbf{1. 评价原理：}将某一指标不同个体测量值按从小到大顺序排列，分为100等份，每1等份代表1个百分位值。适用于正态或非正态分布资料。

\textbf{2. 常用界值点：}$P_3$、$P_{10}$、$P_{25}$、$P_{50}$、$P_{75}$、$P_{90}$、$P_{97}$。

\textbf{3. 等级划分：}

\begin{table}[H]
\centering
\caption{百分位数法五等级划分}
\begin{tabularx}{\textwidth}{ll}
\hline
\textbf{等级} & \textbf{标准} \\
\hline
上等   & $> P_{97}$ \\
中上等 & $P_{75} \sim P_{97}$ \\
中等   & $P_{25} \sim P_{75}$ \\
中下等 & $P_3 \sim P_{25}$ \\
下等   & $< P_3$ \\
\hline
\end{tabularx}
\end{table}

\textbf{4. 优缺点：}

\begin{itemize}[leftmargin=*]
    \item \imp{优点}：适用于任意分布类型的资料；结果直观易懂；国际上广泛使用。
    \item \imp{缺点}：需要较大样本量才能获得稳定的百分位数值。
\end{itemize}
\end{keybox}

\textbf{（三）Z分法（Z-score Method）}

\begin{defbox}
\textbf{Z分法}是一种特殊类型的离差法，以0为中心，用偏离该年龄组标准差的程度反映生长情况。

\textbf{评价原理：}$Z = (X - \bar{x}) / s$，其中$X$为个体实测值，$\bar{x}$为参照人群均值，$s$为参照人群标准差。

\textbf{等级判定（五等级）：}上等$Z>2$；中上等$1<Z\leq2$；中等$-1\leq Z\leq1$；中下等$-2\leq Z<-1$；下等$Z<-2$。

\textbf{优点：}
\begin{itemize}[leftmargin=*]
    \item \imp{跨群体可比}：方便个体、群体间横向比较及同一群体不同年代纵向比较，无需考虑性别、年龄、指标差异。
    \item \imp{分布形态保持}：原始资料转换为Z分后不改变分布形状和数据次序。
    \item \imp{与百分位数对应}：$Z=0$对应$P_{50}$，$Z=1$对应$P_{84.1}$，$Z=2$对应$P_{97.7}$。
\end{itemize}

\textbf{注意事项：}适用于正态分布资料；偏态分布需先进行正态转换（如LMS法）再计算Z分。
\end{defbox}

\subsection{发育速度评价}

\begin{keybox}
\textbf{（一）群体生长速度评价}

\textbf{1. 年增长值（Annual Increment）：}

群体两连续年龄组某指标均值相减，反映该年龄段儿童少年的绝对生长量。

\textbf{2. 年增长率（Annual Growth Rate）：}

年增长值除以该指标的基数（前一年龄组的均值），以百分数表示。计算公式：$V(\%) = (H_{t+1} - H_t) / H_t \times 100\%$。

\textbf{3. 修匀（Smoothing）：}

为消除抽样误差和测量误差对年增长率的影响，对数据进行修匀处理。修匀公式：$V = (a + 2b + c) / 4$，其中 $b$ 为修正后该年龄组年增长率，$a$、$c$ 分别为 $b$ 上下相邻年龄组的年增长率。

\textbf{（二）个体生长速度评价}

常用生长监测图（Growth Monitoring Chart）表示。所用正常值标准应以追踪性资料为基础建立。

评价结果通常以以下方式表示：
\begin{itemize}[leftmargin=*]
    \item \imp{正常}：生长曲线与参考曲线基本平行。
    \item \imp{下降}：生长曲线偏离参考曲线向下。
    \item \imp{缓慢（增长不足）}：生长速率明显低于正常。
    \item \imp{加速}：生长速率明显高于正常。
\end{itemize}
\end{keybox}

\subsection{发育匀称度评价}

\begin{keybox}
\textbf{（一）体型匀称度（Body Type Indices）}

\textbf{1. 身高体重指数 / 克托莱指数（Quetelet Index）：}

$\text{体重(kg)} / \text{身高(cm)} \times 100\%$。反映单位身高的体重，体现人体充实度和营养状况。随年龄增长而增大，男性21岁、女性19岁趋于稳定。

\textbf{2. 身高胸围指数：}

$\text{胸围(cm)} / \text{身高(cm)} \times 100\%$。反映胸廓发育状况和躯干体型。均值在突增高峰前随年龄增长而下降，突增高峰时最低；突增高峰后随年龄增长而上升，成年后趋于稳定。

\textbf{3. 体质量指数（BMI）：}

$\text{体重(kg)} / [\text{身高(m)}]^2$。反映每平方米身体面积包含的体重，敏感反映人体充实度和体型胖瘦，受身高干扰小。广泛用于建立营养不良、超重/肥胖筛查标准。成人 $\geq 24$ 为超重，$\geq 28$ 为肥胖。

\textbf{4. 劳雷尔指数（Rohrer Index）：}

$\text{体重(kg)} / [\text{身高(cm)}]^3 \times 10^7$。反映每$\mathrm{cm}^3$身体的重量，曲线呈"V"字形，能较敏感反映体型胖瘦。缺点是受身材高矮影响大。

\textbf{5. Kaup指数：}

$\text{体重(kg)} / [\text{身长(cm)}]^2 \times 10^4$。主要用于反映婴幼儿体格发育和营养状况，评价儿童机体脂肪总量及分布情况。

\textbf{（二）身材匀称（Body Proportion）}

以坐高/身高比值（SH/H）或躯干/下肢比反映下肢发育状况：

\begin{itemize}[leftmargin=*]
    \item SH/H从婴儿期的0.68逐渐下降至青少年的0.52，提示青春期前下肢较躯干生长快。
    \item SH/H与身高呈显著负相关。
    \item 实际比值 $\leq$ 参照人群值为身材匀称，实际比值 $>$ 参照人群值为不匀称。
    \item 身材匀称评价结果可帮助诊断内分泌及骨骼发育异常疾病。
\end{itemize}
\end{keybox}

\subsection{发育年龄评价}

\begin{keybox}
\textbf{（一）形态年龄（Morphological Age）}

用一般群体某形态指标发育状态制成标准年龄，将个体该形态指标与标准指标比较来评价。最常用的是身高年龄和体重年龄。

\begin{itemize}[leftmargin=*]
    \item \imp{优点}：使用简单，结果明确。
    \item \imp{缺点}：单项形态指标仅反映全身发育一个侧面；单独评价效果常不理想，仅对衡量该体格指标的发育提前或滞后有提示作用。
\end{itemize}
\textbf{（二）齿龄 / 牙齿年龄（Dental Age）}

按儿童牙齿发育顺序制成标准年龄，反映个体发育状况。

\begin{enumerate}[leftmargin=*]
    \item \imp{牙齿萌出数量和质量}：适用于出生后6个月至13岁左右（除第三磨牙外），对早期发育评价有重要价值。
    \item \imp{X线摄片法}：从第1颗牙钙化开始到最后1颗恒牙完成钙化的全过程，划分不同阶段反映牙齿发育水平和速度。该方法准确可靠。
\end{enumerate}

\textbf{（三）骨龄 / 骨骼年龄（Skeletal Age）\difficulty}

骺和小骨骨化中心出现的年龄及骺与骨干愈合的年龄，是反映个体发育水平和成熟程度的精确指标。以手腕部X线摄片最为理想。

\textbf{手腕部骨龄评价的优点：}

\begin{enumerate}[leftmargin=*]
    \item 手、腕骨数目、种类和形状多样，包括长骨、短骨、不规则骨和籽骨，对全身骨骼有很好的代表性。
    \item 手、腕骨各继发性骨化中心出现及掌指骨、尺桡骨的干骺愈合有明显的时间顺序，不同发育阶段间界限明确，易于发现差别。
    \item 拍片方便，投照条件易控制，受检者接受X线剂量小，对保护儿童少年健康有利。
\end{enumerate}

\textbf{常用骨龄评价方法：}

\begin{itemize}[leftmargin=*]
    \item \imp{G-P图谱法（Greulich-Pyle法）}：将手腕部X线片与标准骨龄图谱比较。
    \item \imp{TW系列计分法（Tanner-Whitehouse法）}：对20块手腕骨分别评分后计算总分，再换算为骨龄。TW3为最新版本。
    \item \imp{CHN法（中国人骨龄标准）}：基于中国儿童少年样本建立的骨龄评价标准。
\end{itemize}

\textbf{骨龄的临床应用：}
\begin{itemize}[leftmargin=*]
    \item 评价生长发育成熟度，预测成年身高。
    \item 诊断生长发育障碍（如生长激素缺乏症、性早熟等）。
    \item 指导内分泌疾病的治疗和疗效监测。
    \item 运动员选材和法医学年龄鉴定。
\end{itemize}
\end{keybox}

\subsection{生长发育评价标准的制定方法}

\begin{defbox}
\textbf{常用标准制定方法：}

\begin{enumerate}[leftmargin=*]
    \item \imp{均数标准差法}：适用于正态分布指标，以$\bar{x}\pm s$或$\bar{x}\pm 2s$为界值。

    \item \imp{百分位数法}：适用于正态或偏态分布，以P3、P10、P25、P50、P75、P90、P97等为界值。

    \item \imp{LMS法（偏度-中位数-变异系数法）}：由英国学者Cole提出，是WHO儿童生长标准制定的核心方法。通过L（偏度系数，Box-Cox转换）、M（中位数）、S（变异系数）三个参数，将偏态分布资料转化为正态分布，实现资料的标准化处理，使不同性别、年龄群体的生长发育水平可以进行横向和纵向比较。是目前国际通用的生长发育标准制定方法。
\end{enumerate}
\end{defbox}

% =====================================================================
% 第三部分：健康危险行为监测
% =====================================================================
\section{健康危险行为监测}

\subsection{概述}

\begin{defbox}
\textbf{健康危险行为（Health Risk Behaviors）：}指个体或群体在偏离个人、他人和社会健康期望的方向上存在的一组行为，这些行为会直接或间接地损害青少年的健康、生活质量和生命。

\textbf{健康促进行为（Health Promoting Behaviors）：}指个体或群体表现出的客观上有利于自身和他人健康的行为，包括日常健康行为、保健行为、避免有害环境行为、戒除不良嗜好行为、预警行为等。
\end{defbox}

\begin{keybox}
\textbf{青少年健康危险行为的七种主要类型：}

\begin{enumerate}[leftmargin=*]
    \item \imp{不良饮食行为}：不吃早餐、偏食挑食、过度摄入高热量食物和含糖饮料、不健康减肥行为等。
    \item \imp{缺乏体力活动}：静态生活方式，体育锻炼不足。
    \item \imp{导致伤害的行为}：交通伤害行为（不系安全带、醉酒驾车）、溺水、跌倒、暴力行为、自杀自伤行为等。
    \item \imp{物质滥用行为}：吸烟、饮酒、药物滥用（包括违禁药物和非医疗目的的药物使用）。
    \item \imp{精神成瘾行为}：网络成瘾、游戏成瘾等。
    \item \imp{久坐行为/静态行为}：长时间看电视、使用电脑、玩电子游戏等。
    \item \imp{睡眠不足}：因学业压力、电子产品使用等原因导致的睡眠时间不足和睡眠质量下降。
\end{enumerate}
\end{keybox}

\subsection{监测内容与方法}

\begin{keybox}
\textbf{（一）监测对象与时间}

\begin{itemize}[leftmargin=*]
    \item \imp{监测对象}：以中学生（初中和高中）为主要监测对象，部分指标可延伸至小学高年级和大学生群体。
    \item \imp{监测时间}：一般每1--2年进行一次，也可根据需要加密。通常安排在学校正常教学期间进行。
\end{itemize}

\textbf{（二）监测指标体系}

通常包括七大指标类别：

\begin{enumerate}[leftmargin=*]
    \item \imp{人口学与基本情况}：年龄、性别、年级、家庭背景、学习成绩等。
    \item \imp{饮食行为与营养状况}：膳食频率、早餐习惯、体重控制行为等。
    \item \imp{体力活动与静态行为}：体育锻炼时间、户外活动时间、看电视/玩电子游戏时间等。
    \item \imp{伤害相关行为}：交通安全行为、暴力行为、自杀意念与行为等。
    \item \imp{物质使用行为}：吸烟、饮酒、药物使用等。
    \item \imp{网络与电子媒体使用行为}：上网时间、网络成瘾倾向等。
    \item \imp{睡眠行为}：睡眠时间、睡眠质量、日间嗜睡等。
\end{enumerate}

\textbf{（三）质量控制}

\begin{itemize}[leftmargin=*]
    \item \imp{调查工具标准化}：使用信度和效度经过验证的标准化问卷。
    \item \imp{调查员培训}：统一培训调查员，规范调查用语和程序。
    \item \imp{匿名调查}：保证调查的匿名性，提高回答的真实性。
    \item \imp{数据清理和逻辑检验}：剔除不合格问卷，进行逻辑一致性检查。
\end{itemize}
\end{keybox}

\subsection{监测应用}

\begin{defbox}
\textbf{（一）学生常见病监测}

通过对学生常见病（近视、肥胖、龋齿、脊柱弯曲异常等）的系统监测，掌握学生常见病的流行状况和变化趋势，为制定学校卫生干预策略提供依据。

\textbf{（二）青少年健康危险行为监测系统（Youth Risk Behavior Surveillance System, YRBSS）}

由美国CDC于1990年建立，是国际上最具影响力的青少年健康危险行为监测系统之一。

\begin{itemize}[leftmargin=*]
    \item \imp{监测目的}：监测青少年中导致死亡、残疾和社会问题的优先健康危险行为的流行状况和变化趋势。
    \item \imp{监测内容}：伤害行为、暴力行为、自杀行为、吸烟、饮酒及药物使用、性行为、膳食行为、体力活动。
    \item \imp{监测方法}：每两年一次，采用多阶段整群抽样设计，匿名自填问卷。
    \item \imp{中国实践}：中国自2005年起在全国部分省市开展青少年健康危险行为监测，为制定学校卫生政策和干预措施提供科学依据。
\end{itemize}

\textbf{（三）监测结果的应用}

\begin{enumerate}[leftmargin=*]
    \item 描述青少年健康危险行为的流行特征和变化趋势。
    \item 识别高危人群和重点干预领域。
    \item 评价学校健康教育及干预措施的效果。
    \item 为国家和地方制定青少年健康政策提供数据支持。
    \item 促进"家--校--社"联动的青少年健康促进体系的建立和完善。
\end{enumerate}
\end{defbox}

\newpage
